\begin{proof} [Proof of Theorem~\ref{thm:thin-by-number}]

Observe that $C_{n-1}$ is thin iff $\delta(b) = 1$ for all $b \in
C_{n-1}\backslash C_n$ iff $\delta(b) =
o(C_n)/m(\mathcal{C}, n, b)$.

If $\mathcal{C}$ is thin, then $C_j//C_{j+1}$ is a thin closed subset
for all $0 \leq j \leq n-1$, by definition. So $\delta$ is constant on
each double coset $C_j//C_{j+1}\cdot b//C_{j+1}\cdot C_j//C_{j+1}$ for
all $0 \leq i \leq j \leq n-1$ and all $b \in B\backslash C_i$. Hence, $\delta(b//C_i) =
\frac{\delta(b)}{o(C_i)}m(\mathcal{C}, i, b)$ for all $0 \leq i \leq
n-1$ and $b \in B\backslash C_i$, by Theorem~\ref{thm:constant_2coset}.
But $C_{i-1}//C_i$ thin and $b \in C_{i-1}\backslash C_i$ imply that
$\delta(b//C_i) = 1$. Hence, $\delta(b) = o(C_i)/m(\mathcal{C},
i, b)$ for $0 \leq i \leq n-1$.

Conversely, suppose that $\delta(b) = o(C_i)/m(\mathcal{C}, i, b)$ for
all $0 \leq i \leq n$ and all $b \in C_{i-1}\backslash C_i$. In
particular if $i = n$, then $C_{n-1}$ is thin, as noted above. So
$\delta$ is constant on $C_{n-1}bC_{n-1}$ for all $b \in B$.

If $i < n$ and $b \in C_{i-1}\backslash C_i$, then $C_{n-1}$ thin, Lemma~\ref{lem:constant} and Lemma
\ref{lem:chain_num_mult} imply that
\[
\delta(b//C_{n-1}) \Mk = \Mk  \frac{\delta(b)}{o(C_{n-1})}m(\mathcal{C}, n-1,
b) \Mk =\Mk  \frac{o(C_i)m(\mathcal{C}, n-1, b)}{m(\mathcal{C}, i, b)o(C_{n-1})}
\Mk =\Mk  \frac{o(C_i//C_{n-1})}{m(\mathcal{C}//C_{n-1}, i, b//C_{n-1})}.
\]
So the same hypothesis holds for the chain $\mathcal{C}//C_{n-1}$ in
$B//C_{n-1}$ as for $\mathcal{C}$. Induction on $n$ implies that
$\mathcal{C}//C_{n-1}$ is thin. Since $C_j//C_{j+1} =
(C_j//C_{n-1})//(C_{j+1}//C_{n-1})$ for all $0 \leq j \leq n-2$ by
Proposition~\ref{prop:subset_corresp}\ref{prop2.5_ii}, $\mathcal{C}$ is thin.
\end{proof}
